\section{Conclusions}
\label{sec:conclusions}

A physics-informed digital twin of an industrial top-fired steam methane reformer has been built from published kinetics, plant data
and creep curves, coupling a one-dimensional tube model, a long-furnace model, a Larson--Miller life module, Gaussian-process
surrogates with conformal prediction intervals and Monte Carlo uncertainty propagation in one open, versioned pipeline.

\begin{itemize}
\item The tube model reproduces the industrial case of Xu and Froment within \xfRMSET{} K in gas temperature once the film
  coefficient is matched, and the heat-transfer correlations printed in that paper are found to overpredict the film coefficient by an
  order of magnitude at industrial Reynolds numbers. The coupled model calibrated on three plant operating points reproduces outlet
  temperatures within 4 K and wall temperatures within 12 K, with identifiable parameters and a systematic residual at the upper
  elevation that points to the flame-length response to load.
\item Firing intensity, steam-to-carbon ratio and excess air explain \SI{82}{\percent} of the variance of the hot-spot temperature
  with negligible interactions; inlet pressure enters the life only through the hoop stress.
\item Under outlet-temperature control, load-following and renewable-following cost less than \SI{2}{\percent} of the steady life per
  kmol of hydrogen, short over-firing events at most \SI{3}{\percent}, and steam-to-carbon ratio $\pm 17$\,\%; catalyst ageing
  under a methane-slip target multiplies the life cost by \SxAgeSlip{} in four years.
\item The calibrated base is Pareto-dominated: a knee regime gives \kneeHtwo{}\,\% hydrogen at \kneeFuel{}\,\% fuel per kmol and
  \kneeLife{} of the life cost, and at constant production the life cost can be cut to \isoLife{}. The optimal steam-to-carbon ratio
  depends entirely on the steam credit, moving from 4.0 to 2.7 when steam raising is charged in full.
\item The \SI{90}{\percent} uncertainty on the absolute rupture time spans a factor of 50, dominated by creep-curve scatter
  (\shareCreep{}\,\%), while paired ratios between regimes remain tight. The ageing, knee and steam-to-carbon conclusions hold in
  \SI{81}{\percent} to \SI{91}{\percent} of the samples; the flexible-operation conclusion does not survive a load-dependent flame
  length.
\end{itemize}

The next steps are the transcription of the NIMS rupture data for the HP-Nb alloys into the prepared ingestion path, a per-tube
extension with a measured wall-temperature distribution, and, above all, a multi-elevation pyrometer survey at several firing
rates to identify the flame-length response that decides the life cost of flexible operation.
